\section{Conclusion}


In this paper, we proposed a unified dual-mask physical model to address the severe performance degradation of Epipolar Plane Image (EPI) based light field depth estimation caused by non-Lambertian surfaces and complex textures. By decoupling material properties from geometric structures, our approach effectively mitigates the failure of the Lambertian reflectance assumption in mixed scenarios. Extensive evaluations demonstrate that the proposed framework, equipped with domain-balanced sampling, achieves an overall Mean Absolute Error (MAE) of 0.133 and a mixed urban scene MAE of 0.081, validating its efficacy in maintaining robust depth estimation across heterogeneous environments.

The specific contributions of this work are summarized as follows:
1) \textbf{Three-layer angular signal decomposition and geometric-material decoupling:} We established a theoretical framework that reveals the intrinsic structure of light field angular signals from physical and geometric perspectives. By substituting frequency-domain features with geometric features, we overcame the coarseness of frequency representations under low angular resolutions, enabling efficient and highly accurate material classification to support the dual-branch network design.
2) \textbf{Dual-mask and geometry-routed architecture for mixed-scene depth estimation:} We developed a hybrid estimation framework that alleviates the performance drop of conventional single-EPI networks in mixed-material scenes. Through explicit mask learning and feature routing, the model preserves the advantage of EPI slope extraction in Lambertian regions while applying differentiated processing to non-Lambertian areas, significantly improving depth accuracy and robustness in complex urban environments.
3) \textbf{Root cause analysis and physical boundary definition of EPI assumption failure:} We provided a systematic delineation of the physical limits of EPI-based methods, supported by empirical evidence showing MAEs of 0.387 and 0.411 for complex Lambertian and non-Lambertian regions, respectively. This analysis confirms that purely modifying network architectures or loss functions cannot overcome the physical bottlenecks of specular highlights and scattering, redirecting future research efforts away from ineffective computational investments in pure EPI frameworks.

Moving forward, advancing light field depth estimation beyond these physical boundaries will necessitate the integration of non-EPI paradigms, such as neural radiance fields or learned feature matching, coupled with large-scale non-Lambertian datasets to achieve robust 3D perception in unconstrained environments.


The transition from frequency-domain analysis to spatial-domain geometric features represents a fundamental shift in how light field (LF) systems process angular information. Traditional methods relying on 2D Discrete Fourier Transforms implicitly demand dense angular sampling to satisfy the Nyquist-Shannon sampling theorem, a requirement physically unattainable for compact plenoptic cameras constrained to low angular resolutions (e.g., $9 \times 9$). By formulating the three-layer angular signal decomposition as:
\begin{equation}
\label{eq:eq7}
I(u,v) = DC + a\cdot u + b\cdot v + \epsilon(u,v)
\end{equation}
our framework circumvents these hardware limitations. The residual component $\epsilon(u,v)$, when analyzed via spatial geometric features like angular gradients, provides a computationally lightweight yet physically grounded mechanism to isolate material reflectance. This decoupling not only advances algorithmic accuracy but also broadens the deployability of LF depth estimation on resource-constrained edge devices where high-density microlens arrays are impractical.

From a systems integration perspective, the proposed geometry-routed dual-branch architecture introduces conditional computation into the depth estimation pipeline. While dynamically routing features based on material masks significantly enhances performance in heterogeneous environments, it inherently induces branch divergence. In strict real-time applications such as autonomous driving or robotic navigation, variable execution paths can lead to unpredictable latency and suboptimal GPU utilization. Future system-level optimizations must address this by employing static graph compilation techniques or batch-padding strategies to ensure deterministic inference times without sacrificing the adaptive benefits of the dual-mask routing.

Despite robust empirical performance, the unified physical model exhibits certain boundary limitations. The linear superposition assumption inherent in the three-layer decomposition effectively models direct illumination and primary specular reflections (e.g., the characteristic X-pattern in EPIs). However, it struggles with complex multi-bounce optical phenomena, such as caustics, deep translucent subsurface scattering, or extreme interreflections. In these scenarios, the residual term $\epsilon(u,v)$ becomes highly non-linear and spatially entangled with the geometric disparity component $a\cdot u + b\cdot v$, degrading the efficacy of the material classifier. Furthermore, severe occlusions disrupt the continuous epipolar geometry, corrupting the angular correlation features required for accurate mask generation.

Looking forward, the geometric-material decoupling paradigm established in this work opens several promising research avenues. A natural extension involves incorporating the temporal dimension to handle dynamic light fields, where motion blur and temporal occlusions further complicate the EPI structure. Additionally, the explicit material masks and decoupled geometric priors generated by our model could serve as powerful physical regularizers for modern implicit and explicit 3D reconstruction frameworks, such as Neural Radiance Fields (NeRFs) and 3D Gaussian Splatting (3DGS). By injecting these physically derived material boundaries into the rendering equations, future systems could achieve unprecedented fidelity in novel view synthesis and scene reconstruction for complex, mixed-material real-world environments.

Despite the demonstrated efficacy of the unified dual-mask physical model in decoupling material reflectance from geometric disparity, several limitations remain that outline promising directions for future research. 

First, while the proposed three-layer angular signal decomposition model $I(u,v) = DC + a\cdot u + b\cdot v + \varepsilon(u,v)$ effectively circumvents the coarseness of frequency-domain features at standard low angular resolutions (e.g., $9 \times 9$), its reliance on spatial-domain geometric features (such as angular gradients and coefficients of variation) becomes vulnerable under extremely sparse angular sampling (e.g., $4 \times 4$ or $2 \times 2$). In such ultra-sparse regimes, the discrete approximation of angular derivatives degrades, potentially compromising the accuracy of the material-aware masks. Future work will investigate adaptive feature extraction mechanisms or angular super-resolution priors to maintain robust physical decoupling under severely restricted angular bandwidths.

Second, the current formulation implicitly assumes that the residual component $\varepsilon(u,v)$ primarily encapsulates non-Lambertian material properties and high-frequency sensor noise. However, in the presence of complex global illumination phenomena such as multi-bounce interreflections, caustics, or refractive translucent materials, $\varepsilon(u,v)$ may contain structured geometric cues that violate the local linearity assumption of the Epipolar Plane Image. Consequently, the geometry-routed architecture might misclassify these regions or yield suboptimal disparity estimates due to conflicting structural signals. Extending the physical model to incorporate higher-order light transport equations or adopting a data-driven hierarchical residual decomposition could further disentangle these complex optical effects from pure geometric disparity, thereby enhancing robustness in highly reflective or transparent environments.

Finally, the current framework processes light field data on a strictly per-frame basis, overlooking the rich temporal correlations inherent in multi-view video systems and dynamic autonomous driving scenarios. Non-Lambertian artifacts, such as dynamic specular highlights, exhibit distinct temporal trajectories and parallax behaviors compared to diffuse surface points when the camera platform or scene objects are in motion. Integrating the proposed dual-mask architecture with spatiotemporal consistency constraints represents a critical next step for video-oriented applications. By formulating a 5D plenoptic function model that incorporates time, future iterations can leverage temporal epipolar geometry and optical flow priors to further regularize depth estimation in dynamic, mixed-material environments. This extension would ultimately advance the deployment of robust, real-time 3D perception systems in complex, real-world autonomous navigation and virtual reality applications.